\section{Scalar Products}

\begin{definition}[Index Notation]
  \(A\) vector equation such as

\[
\boldsymbol{c}=\boldsymbol{a}+\boldsymbol{b}
\]

is written in index notation as

\[
c_i=a_i+b_i,
\]

where it is understood that the equation holds for \(i=1, \ldots, n\), meaning an arbitrary component in the standard (Cartesian) basis \(\left\{\boldsymbol{e}_1, \ldots, \boldsymbol{e}_n\right\}\).
\end{definition}

\begin{remark}
  An index that appears once in each term is called a free index. The choice of letter is arbitrary, but all terms in an equation must have the same free index (or indices).
\end{remark}

\begin{definition}[Einstein Summation Convention]
  In index notation, we write the scalar product as

\[
\boldsymbol{a} \cdot \boldsymbol{b}=a_j b_j,
\]

where the repeated index indicates that the term should be summed from \(j=1\) to \(n\). This implied summation is called the \textbf{(Einstein) summation convention}.

An index that appears exactly twice in a term is called a \textbf{dummy index}, and always implies a summation over that index. Again, the choice of letter is arbitrary.
\end{definition}

\begin{definition}[Kronecker Delta]
  \[
\boldsymbol{a} \cdot \boldsymbol{b}=\left(a_1, a_2, a_3\right)\left(\begin{array}{lll}
1 & 0 & 0 \\
0 & 1 & 0 \\
0 & 0 & 1
\end{array}\right)\left(\begin{array}{l}
b_1 \\
b_2 \\
b_3
\end{array}\right)=a_1 b_1+a_2 b_2+a_3 b_3 .
\]

In index notation, this may be written

\[
\boldsymbol{a} \cdot \boldsymbol{b}=\delta_{i j} a_i b_j,
\]

where the components of the identity matrix are given by the Kronecker delta, defined as

\[
\delta_{i j}= \begin{cases}1 & \text { if } i=j, \\ 0 & \text { if } i \neq j .\end{cases}
\]
\end{definition}

\begin{proposition}
  Multiplying an expression with free index \(i\) by \(\delta_{i j}\) is equivalent to replacing \(i \rightarrow j\). We can simplify expressions by replacing the Kronecker Delta with a replacement instead.
\end{proposition}

\begin{proofidea}
  
  
\end{proofidea}
  
\begin{formalproof}
  
\end{formalproof}

\section{Vector Products} 

\begin{definition}[Alternating Tensor (Levi-Civita) Symbol]
  In \(\mathbb{R}^3\), we can write vector products in index notation by introducing the alternating tensor (or Levi-Civita symbol),

\[
\epsilon_{i j k}= \begin{cases}0 & \text { if any of } i, j, k \text { are equal, } \\ +1 & \text { if }(i, j, k)=(1,2,3),(2,3,1) \text { or }(3,1,2), \\ -1 & \text { if }(i, j, k)=(1,3,2),(2,1,3) \text { or }(3,2,1) .\end{cases}
\]
\end{definition}

\begin{remark}
  
This allows us to express the vector product in index notation as

\[
[\boldsymbol{a} \times \boldsymbol{b}]_i=\epsilon_{i j k} a_j b_k .
\]

\end{remark}

\begin{example}
  Note that this scalar triple product is a scalar, so there should be no free index.
We have

\[
\begin{aligned}
\boldsymbol{a} \cdot(\boldsymbol{b} \times \boldsymbol{c}) & =a_i[\boldsymbol{b} \times \boldsymbol{c}]_i \\
& =a_i \epsilon_{i j k} b_j c_k \\
& =\epsilon_{i j k} a_i b_j c_k \quad[\text { order doesn't matter for scalars }] \\
& =\epsilon_{k i j} a_i b_j c_k \\
& =c_k \epsilon_{k i j} a_i b_j \\
& =c_i \epsilon_{i j k} a_j b_k \quad[\text { relabelling the dummy indices }(i, j, k) \rightarrow(j, k, i)] \\
& =c_i[\boldsymbol{a} \times \boldsymbol{b}]_i \\
& =\boldsymbol{c} \cdot(\boldsymbol{a} \times \boldsymbol{b}) .
\end{aligned}
\]

The final relabelling wasn't strictly necessary, just aesthetically pleasing!
\end{example}

\begin{remark}
  You can relabel dummy indicies as needed to aid you.
\end{remark}

\begin{proposition}[The \(\epsilon - \delta\) identity]
  \[
\epsilon_{i j k} \epsilon_{k l m}=\delta_{i l} \delta_{j m}-\delta_{i m} \delta_{j l} .
\]
\end{proposition}

\begin{proofidea}
  
  
\end{proofidea}
  
\begin{formalproof}
  
\end{formalproof}

\section{Derivatives}

\begin{definition}[Grad, Div and Curl in Index Notation]
  \[
\boldsymbol{\nabla} f=\frac{\partial f}{\partial x} \boldsymbol{e}_1+\frac{\partial f}{\partial y} \boldsymbol{e}_2+\frac{\partial f}{\partial z} \boldsymbol{e}_3, \quad \boldsymbol{\nabla} \cdot \boldsymbol{f}=\frac{\partial f_1}{\partial x}+\frac{\partial f_2}{\partial y}+\frac{\partial f_3}{\partial z}, \quad \boldsymbol{\nabla} \times \boldsymbol{f}=\left|\begin{array}{ccc}
\boldsymbol{e}_1 & \boldsymbol{e}_2 & \boldsymbol{e}_3 \\
\frac{\partial}{\partial x} & \frac{\partial}{\partial y} & \frac{\partial}{\partial z} \\
f_1 & f_2 & f_3
\end{array}\right|
\]

Therefore the gradient, divergence and curl are written in index notation as

\[
[\nabla f]_i=\frac{\partial f}{\partial x_i}, \quad \nabla \cdot f=\frac{\partial f_j}{\partial x_j}, \quad[\nabla \times f]_i=\epsilon_{i j k} \frac{\partial f_k}{\partial x_j} .
\]
\end{definition}

\begin{remark}
  We can complete the product rule proof here using this notation

The \(i\) component of the left-hand side is

\[
\begin{aligned}
{[\boldsymbol{\nabla} \times(\boldsymbol{f} \times \boldsymbol{g})]_i } & =\epsilon_{i j k} \frac{\partial}{\partial x_j}\left(\epsilon_{k l m} f_l g_m\right) \\
& =\epsilon_{i j k} \epsilon_{k l m} \frac{\partial}{\partial x_j}\left(f_l g_m\right) \\
& =\left(\delta_{i l} \delta_{j m}-\delta_{i m} \delta_{j l}\right) \frac{\partial}{\partial x_j}\left(f_l g_m\right) \quad[\text { by Proposition 5.2 }] \\
& =\left(\delta_{i l} \delta_{j m}-\delta_{i m} \delta_{j l}\right)\left(\frac{\partial f_l}{\partial x_j} g_m+f_l \frac{\partial g_m}{\partial x_j}\right) \\
& =\delta_{i l} \delta_{j m}\left(\frac{\partial f_l}{\partial x_j} g_m+f_l \frac{\partial g_m}{\partial x_j}\right)-\delta_{i m} \delta_{j l}\left(\frac{\partial f_l}{\partial x_j} g_m+f_l \frac{\partial g_m}{\partial x_j}\right) \\
& =\frac{\partial f_i}{\partial x_j} g_j+f_i \frac{\partial g_j}{\partial x_j}-\frac{\partial f_j}{\partial x_j} g_i-f_j \frac{\partial g_i}{\partial x_j} \quad[\text { by Proposition 5.1 }] \\
& =\frac{\partial f_i}{\partial x_j} g_j+f_i(\boldsymbol{\nabla} \cdot \boldsymbol{g})-(\boldsymbol{\nabla} \cdot \boldsymbol{f}) g_i-f_j \frac{\partial g_i}{\partial x_j} .
\end{aligned}
\]
This time the right-hand side contains "triple products", so we start with that and simplify. We have

\[
\begin{aligned}
{[\boldsymbol{f} \times(\boldsymbol{\nabla} \times \boldsymbol{g})]_i } & =\epsilon_{i j k} f_j \epsilon_{k l m} \frac{\partial g_m}{\partial x_l} \\
& =\epsilon_{i j k} \epsilon_{k l m} f_j \frac{\partial g_m}{\partial x_l} \\
& =\left(\delta_{i l} \delta_{j m}-\delta_{i m} \delta_{j l}\right) f_j \frac{\partial g_m}{\partial x_l} \\
& =\delta_{i l} \delta_{j m} f_j \frac{\partial g_m}{\partial x_l}-\delta_{i m} \delta_{j l} f_j \frac{\partial g_m}{\partial x_l} \quad \text { [by Proposition 5.2] } \\
& =f_j \frac{\partial g_j}{\partial x_i}-f_j \frac{\partial g_i}{\partial x_j} \quad \text { [by Proposition 5.1]. }
\end{aligned}
\]

So the \(i\) component of the full right-hand side of \((v i)\) is

\[
\begin{aligned}
{\left[\mathrm{RHS}_i\right.} & =g_j \frac{\partial f_i}{\partial x_j}+f_j \frac{\partial g_i}{\partial x_j}+f_j \frac{\partial g_j}{\partial x_i}-f_j \frac{\partial g_i}{\partial x_j}+g_j \frac{\partial f_j}{\partial x_i}-g_j \frac{\partial f_i}{\partial x_j} \\
& =f_j \frac{\partial g_j}{\partial x_i}+g_j \frac{\partial f_j}{\partial x_i} \\
& =\frac{\partial}{\partial x_i}\left(f_j g_j\right)=[\boldsymbol{\nabla}(\boldsymbol{f} \cdot \boldsymbol{g})]_i .
\end{aligned}
\]

\end{remark}


\section{Second Derivatives}

\begin{definition}[Second Derivatives]
  For a scalar field \(g: \mathbb{R}^3 \rightarrow \mathbb{R}\) or a vector field \(\boldsymbol{f}: \mathbb{R}^3 \rightarrow \mathbb{R}^3\), there are five ways to combine two derivatives to get a second derivative operator:

\[
\boldsymbol{\nabla} \cdot(\boldsymbol{\nabla} g), \quad \boldsymbol{\nabla} \times(\boldsymbol{\nabla} g), \quad \boldsymbol{\nabla}(\boldsymbol{\nabla} \cdot \boldsymbol{f}), \quad \boldsymbol{\nabla} \cdot(\boldsymbol{\nabla} \times \boldsymbol{f}), \quad \text { or } \quad \boldsymbol{\nabla} \times(\boldsymbol{\nabla} \times \boldsymbol{f}) .
\]

However, we already know that \(\boldsymbol{\nabla} \times(\boldsymbol{\nabla} g)=\mathbf{0}\) and \(\boldsymbol{\nabla} \cdot(\boldsymbol{\nabla} \times \boldsymbol{f})=0\), so there are only three non-trivial operators:

\begin{itemize}
  \item The Laplacian of a scalar field is the operator defined by

\[
\nabla^2 g=\nabla \cdot(\nabla g) .
\]
\item Next, consider the operator \(\boldsymbol{\nabla}(\boldsymbol{\nabla} \cdot \boldsymbol{f})\). Its \(i\) component is simply

\[
[\boldsymbol{\nabla}(\boldsymbol{\nabla} \cdot \boldsymbol{f})]_i=\frac{\partial}{\partial x_i}\left(\frac{\partial f_j}{\partial x_j}\right) .
\]
\item Finally we consider:

  \[\nabla \times(\nabla \times f)=\nabla(\nabla \cdot f)-\nabla^2 f .\]

\end{itemize}

\end{definition}

\begin{proposition}
  In orthogonal curvilinear coordinates \((u, v, w)\), the Laplacian is given by

\[
\nabla^2 g=\frac{1}{h_u h_v h_w}\left[\frac{\partial}{\partial u}\left(\frac{h_v h_w}{h_u} \frac{\partial g}{\partial u}\right)+\frac{\partial}{\partial v}\left(\frac{h_u h_w}{h_v} \frac{\partial g}{\partial v}\right)+\frac{\partial}{\partial w}\left(\frac{h_u h_v}{h_w} \frac{\partial g}{\partial w}\right)\right] .
\]
\end{proposition}

\begin{proofidea}
  
  
\end{proofidea}
  
\begin{formalproof}
  
\end{formalproof}

\begin{remark}
  Be very careful with \(\boldsymbol{\nabla}^2 \boldsymbol{f}\) in non-Cartesian coordinate systems! For example, in orthogonal curvilinear coordinates, \(\left[\nabla^2 f\right] \cdot e_u \neq \nabla^2 f_u\). To find the correct expression, you would rearrange and calculate

\[
\nabla^2 f=\nabla(\nabla \cdot f)-\nabla \times(\nabla \times f) .
\]

This shows, incidentally, that \(\boldsymbol{\nabla}^2 \boldsymbol{f}\) really is coordinate independent, since everything on the righthand side is.
\end{remark}

